
\begin{table*}[!h]
\centering
\small
\caption{Construct representations compared across validation settings. FeedbackQA evaluates \textit{Answer Quality}; Barter Deals evaluates \textit{Deal Attractiveness}. The table summarizes how each target construct is specified in the criterion prompt and how many LLM-derived scores are produced per artefact.}
\label{tab:construct_representations}
\begin{tabularx}{\textwidth}{@{} 
    l 
    >{\raggedright\arraybackslash}p{3.8cm} 
    >{\raggedright\arraybackslash}X 
    c 
    @{}}
\toprule
\textbf{Dataset} & \textbf{Representation} & \textbf{Criterion content} & \begin{tabular}{@{}c@{}}\textbf{Scores /} \\ \textbf{artefact}\end{tabular} \\
\midrule
\textbf{FeedbackQA} 
& Naive holistic 
& Overall answer quality 
& 1 \\
\addlinespace
& Informed holistic control 
& Alignment; Coverage; Concentration; Clarity 
& 1 \\
\addlinespace
& Formative 
& Alignment; Coverage; Concentration; Clarity 
& 4 \\

\midrule

\textbf{Barter Deals} 
& Naive holistic 
& Overall deal attractiveness 
& 1 \\
\addlinespace
& Macro-formative 
& Reward; Cost; Pitch Quality 
& 3 \\
\addlinespace
& Fine-grained formative 
& Ten dimensions nested within Reward, Cost, and Pitch Quality 
& 10 \\
\bottomrule
\end{tabularx}

\vspace{0.15cm}
\begin{flushleft}
\footnotesize
\textit{Note.} Each score is obtained from one LLM call. Holistic representations produce one overall score, while formative representations produce one score per criterion. Criterion-level scores are recomposed through regression when assessing convergent or predictive validity.
\end{flushleft}
\end{table*}
